\begin{table}[H]
\centering
\scriptsize
\caption{Comparativo entre cenários de rede (regime determinístico).}
\label{tab:resumo_cenarios}
\begin{tabular}{|l|c|c|c|}
\hline
\textbf{Tecnologia} & \textbf{Ideal: média / p95 (ms) / sucesso} & \textbf{Realista: média / p95 (ms) / sucesso} & \textbf{Crítico: média / p95 (ms) / sucesso} \\ \hline
WebSocket & 0,55 / 0,64 / 100,0\% & 86,44 / 99,34 / 100,0\% & 506,52 / 1.386,20 / 99,7\% \\ \hline
gRPC & 0,73 / 0,82 / 100,0\% & 84,97 / 100,86 / 100,0\% & 525,69 / 1.329,79 / 100,0\% \\ \hline
SOAP & 5,96 / 23,51 / 100,0\% & 205,41 / 514,36 / 100,0\% & 1.144,59 / 2.721,22 / 99,7\% \\ \hline
REST & 6,03 / 23,08 / 100,0\% & 192,38 / 293,17 / 100,0\% & 1.150,82 / 2.638,85 / 99,7\% \\ \hline
SSE & 6,70 / 24,99 / 100,0\% & 211,25 / 513,55 / 100,0\% & 1.145,98 / 2.778,87 / 99,7\% \\ \hline
GraphQL & 7,93 / 23,46 / 100,0\% & 189,59 / 205,74 / 100,0\% & 1.096,56 / 2.300,39 / 100,0\% \\ \hline
Webhook & 15,83 / 59,35 / 100,0\% & 181,80 / 195,16 / 100,0\% & 1.150,38 / 2.298,24 / 99,7\% \\ \hline
\end{tabular}
{\footnotesize\\ \textbf{Fonte: Dados da pesquisa (2026).}}
\end{table}
